\section{Q-Learning}
    In this section we learn another TD-learning algorithm, the \textbf{Q-Learning}.
    Unlike Sarsa, Q-Learning bootstraps the q value using the maximum action value of current state.
    \subsection{Main Idea}
        The Q-Learning uses following update equation for action values:
        \begin{align}
            q_{t+1}(s_t,a_t) &= q_t(s_t,a_t) + \alpha_t(s_t,a_t)\Big[ r_{t+1} + \gamma \max_{a \in \mathcal{A}(s_{t+1},a)}{q_t(s_{t+1},a)} - q_t(s_t,a_t) \Big] \\
            q_{t+1}(s,a) &= q_t(s,a),\text{ for all }(s,a) \neq (s_t,a_t)
        \end{align}

        The Q-Learning solves following Bellman optimality equation:
        \begin{equation}
            q(s,a) = \E\Big[ R_{t+1} + \gamma \max_aq(S_{t+1},a)|S_t=s,A_t=a \Big]
        \end{equation}
        \begin{proof}
            Check for \cite{Zhao2025} pg.140.
        \end{proof}

        \begin{rem}{On-policy and Off-policy}{onPoloffPol}
            As one can see, the Sarsa solves Bellman equation and then make policy update, but Q-Learning directly learns the optimal policy.
            The \textbf{behavior policy}, is used to make the trajectory. The \textbf{target policy} is the policy that learns the optimal policy.
            When the behavior policy and target policy equals, the algorithms is called \textbf{on-policy} algorithm. If not, it is called \textbf{off-policy}.
            Off policy can separate the exploration with target policy. The target policy only has to care about choosing the best action in current situation.
            Since the off-policy uses policy other than its target policy, it requires some techniques such as \textbf{importance sampling} and \textbf{tree-backup} algorithms.
            These will be covered later.
            
            The on-policy algorithms cannot be implemented as an \textbf{offline} fashion, since it requires real-time interaction with environment.
            The off-policy algorithms, in contrast, can be implemented as an offline or online fashion.
        \end{rem}

    \subsection{Pseudocode}
        \Algorithm{qlearning}
        {Q-Learning}
        {
            \inputitem{$\alpha_t(s,a) \in (0,1]$}{step size parameter satisfying RM condition} \\
            \inputitem{$b$}{behavior policy} \\
            \inputitem{$\pi$}{target policy}
        }
        {
            \inputitem{$q_*$}{the optimal policy}
        }
        {
            \tcp{initialize}
            \ForEach{$(s,a) \in \mathcal{S} \times \mathcal{A}$}
            {$q(s,a) \leftarrow
            \begin{cases}
                \text{arbitrary,} &\text{ if $s$ is not in terminal states} \\
                0 &\text{ if $s$ is in terminal states}
            \end{cases}$\;}
            $s \leftarrow$ initialize\;
            \ForEach{episode} {
                \Repeat{$s$ is not terminal state} {
                    choose action $a$ from $s$ using behavior policy $b$\;
                    observe state $s'$ and reward $r$\;
                    \tcp{update q value}
                    $q(s,a) \leftarrow q(s,a) + \alpha_t(s,a)\Big[ r + \gamma \max_a{q(s',a)} - q(s,a) \Big]$
                    \tcp{update target policy}
                    update target policy according to $q$\;
                    $s \leftarrow s'$\;
                }
            }
        }
    \subsection{Implementation Detail}
        The behavior policy can be any policy that explores well, but in practice we use decaying $\varepsilon$-greedy due to the convergence rate. (I guess...)
    \subsection{Convergence}
        TODO: Convergence of Q-Learning
    \subsection{Pros and Cons}
        Pros:
        \begin{itemize}
            \item It is off-policy. Whatever the behavior policy does, the target policy learns the optimal policy.
            \item Does not fears getting penalty, as opposed to Sarsa. This can be checked in \cite{Sutton2020}, cliff walking.
        \end{itemize}
        Cons:
        \begin{itemize}
            \item Can suffer from \textbf{maximization bias}, since the Q-Learning uses maximum q value for bootstrapping. That is, the agent keeps
                overestimate the q values, which can lead to slow convergence. For more information, see \cite{Sutton2020}. To cope with this, we can use \textbf{Double Q-Learning}.
        \end{itemize}